%---------------------------------------------------------------------
% Resume Template Author: Anubhav Singh | License: MIT
% Resume Author: Siqi Guo — Professor Outreach Version, Feb 2026
%---------------------------------------------------------------------

% ================ Document Configuration ==============
\documentclass[a4paper,20pt]{article}

\usepackage{geometry}
\geometry{left=1.2cm, top=1.2cm, right=1cm, bottom=1cm}

\usepackage{latexsym}
\usepackage[empty]{fullpage}
\usepackage{titlesec}
\usepackage{marvosym}
\usepackage[usenames, dvipsnames]{color}
\usepackage{verbatim}
\usepackage{enumitem}
\usepackage{hyperref}
\usepackage{fancyhdr}
\usepackage{subfig, graphicx}
\usepackage{fontawesome5}
\usepackage{setspace}
\usepackage{amsmath}
\usepackage{multicol}

\pagestyle{fancy}
\fancyhf{}
\fancyfoot{}
\renewcommand{\headrulewidth}{0pt}
\renewcommand{\footrulewidth}{0pt}

% Adjust margins (commented out as they reduce margins)
\addtolength{\oddsidemargin}{-0.530in}
\addtolength{\evensidemargin}{-0.375in}
\addtolength{\textwidth}{1in}
\addtolength{\topmargin}{-.45in}
\addtolength{\textheight}{1in}

\urlstyle{rm}
\raggedbottom
\raggedright
\setlength{\tabcolsep}{0in}

%================ Customized commands =================

\renewcommand{\labelitemii}{$\circ$}

\titleformat{\section}
    {\vspace{-10pt} \scshape \raggedright \large}
    {}
    {0em}
    {}
    [\color{black} \titlerule \vspace{-6pt}]

% region------------ Education SubHeading ----------
\newcommand{\resumeEducationSubHeadingListStart}{
  \begin{description}
}
\newcommand{\resumeEducationSubHeadingListEnd}{
  \end{description}
}
\newcommand{\resumeEducationSubheading}[4]{%
  \vspace{-1pt}\item
    \begin{tabular*} {0.98\textwidth} {@{\hspace{-0.5em}}p{0.8\textwidth}@{\extracolsep{\fill}}r@{}}
        \textbf{#1} & {#2} \\
               {#3} & {#4} \\
    \end{tabular*}%
  \vspace{-5pt}%
}
\newcommand{\courseList}[1]{%
  \vspace{0pt}
  \begin{description}[leftmargin=-25pt, style=unboxed]
    Relevant Courses: \small{#1} 
  \end{description}
  \vspace{-5pt}
}
% endregion

% region----------- Publications SubHeading ----------
\newcommand{\resumePublicationsSubHeadingListStart}{
  \begin{itemize} [label={}, leftmargin=*, topsep=0pt, itemsep=2pt, parsep=0pt]
}
\newcommand{\resumePublicationsSubHeadingListEnd}{
  \end{itemize}
}
\newcommand{\resumePublicationSubItem}[2]{%
  \resumePublicationItem{#1}{#2}
}
\newcommand{\resumePublicationItem}[2]{
  \item\hspace{-0.8em}{
    {#1}{\quad {#2}}
  }
}
% endregion

% -------------- Skills SubHeading ------------
\newcommand{\resumeSkillsSubHeadingListStart}{
  \begin{description}
}
\newcommand{\resumeSkillsSubHeadingListEnd}{
  \end{description}
}
\newcommand{\resumeSkillsSubItem}[2]{%
  \resumeSkillsItem{#1}{#2}
  \vspace{-3pt}%
}
\newcommand{\resumeSkillsItem}[2]{
  \item\hspace{-0.8em}{
    \textbf{#1}{: \small{#2} \vspace{-2pt}}
  }
}

% ------------ Research SubHeading ------------
\newcommand{\resumeResearchSubHeadingListStart}{
  \begin{itemize} [label={}, leftmargin=*]
}
\newcommand{\resumeResearchSubHeadingListEnd}{
  \end{itemize}
}
\newcommand{\resumeResearchSubheading}[6]{
  \vspace{5pt}\item
    \begin{tabular*} {0.98\textwidth} {@{\hspace{-0.5em}}l@{\extracolsep{\fill}}r@{}}
      \textbf{#1} & {#2}\null \\
      {#3} & {#4} \\
      {#5} & {#6}
    \end{tabular*}
  \vspace{0pt}
}
% 4-arg variant: no advisor/institution row
\newcommand{\resumeResearchSubheadingShort}[4]{
  \vspace{5pt}\item
    \begin{tabular*} {0.98\textwidth} {@{\hspace{-0.5em}}l@{\extracolsep{\fill}}r@{}}
      \textbf{#1} & {#2}\null \\
      {#3} & {#4}
    \end{tabular*}
  \vspace{-6pt}
}
\newcommand{\resumeResearchItemListStart}{
  \begin{itemize} [leftmargin=*]
}
\newcommand{\resumeResearchItemListEnd}{
  \end{itemize}
  \vspace{-5pt}
}
\newcommand{\resumeResearchItem}[1]{
  \item{
    {#1 \vspace{-1pt}}
  }
}
\newcommand{\resumeResearchMidItemListStart}{
  \begin{itemize} [leftmargin=*]
}
\newcommand{\resumeResearchMidItemListEnd}{
  \end{itemize}
  \vspace{-5pt}
}
\newcommand{\resumeResearchMidItem}[2]{
  \item{
    \textbf{#1}
    {#2 \vspace{-2pt}}
  }
}

\setlist[itemize, 2]{label=$\bullet$}
% ------------ Projects SubHeading List -------------
\newcommand{\resumeProjectsSubHeadingListStart}{
  \begin{itemize} [leftmargin=*, label={}]
}
\newcommand{\resumeProjectsSubHeadingListEnd}{
  \end{itemize}
}
\newcommand{\resumeProjectsSubHeading}[2]{
  \resumeProjectsSubHeadingItem{#1}{#2}
  \vspace{-10pt}
}
\newcommand{\resumeProjectsSubHeadingItem}[2]{
  \vspace{-1pt}
  \item{
    \begin{tabular*} {0.98\textwidth} {@{\hspace{-0.5em}}l@{\extracolsep{\fill}}r}
      \textbf{#1} & {#2} \\
    \end{tabular*}
  \vspace{-8pt}
  }
}
\newcommand{\resumeProjectsItemListStart}{
  \begin{itemize} [leftmargin=*]
}
\newcommand{\resumeProjectsItemListEnd}{
  \end{itemize}
  \vspace{0pt}
}
\newcommand{\resumeProjectsItem}[1]{
  \item{
    {#1} \vspace{-2pt}
  }
}

%=================================================
%%%%%%%%%%%  CV (Resume) STARTS HERE  %%%%%%%%%%%%
\begin{document}

%-----------------HEADING-----------------
\begin{tabular*} {\textwidth} {l@{\extracolsep{\fill}}r}
  \textbf{ \centerline {\LARGE Siqi  Guo} \vspace {5pt} } \\
  \centerline {\faEnvelope ~: \href{mailto:siqiguo@andrew.cmu.edu}{siqiguo@andrew.cmu.edu} | \faIcon{linkedin} \href{https://www.linkedin.com/in/siqiguo047/?locale=en_US}{: linkedin.com/in/siqiguo047} | \faPhone ~: (+1) 412-612-0495 \quad | \faGlobe \href{https://guo-lab.github.io}{: https://guo-lab.github.io} } \\
\end{tabular*}

\vspace{0.2em}
\begin{center}
  \parbox{0.96\textwidth}{
    \textbf{Summary:}
    CMU M.S. ECE graduate with expertise in embedded control, bio-inspired locomotion, and autonomous planning. Currently developing parallelized \textbf{C++ IK and motion planning pipelines} for wheeled humanoid warehouse robots at Yondu AI. Previously embedded software engineer at Cepton (LiDAR systems) and research assistant in the CMU Biorobotics Lab (\textbf{EigenBot}, \textbf{MMPUG}). Published in ICRA 2025, AAAI 2023, and IEEE Neural Networks. Seeking research collaboration in Embodied AI systems, world models for robotics, real-time robot learning, and scalable deployment of foundation models on physical platforms.
  }
\end{center}
\vspace{-1.5em}

\vspace{1em}

%----------------EDUCATION-----------------
\section{Education}
\resumeEducationSubHeadingListStart
\resumeEducationSubheading
{Carnegie Mellon University} {Pittsburgh, PA}
{Master of Science in Electrical and Computer Engineering; \quad GPA: 3.91/4.0} {May 2025}
\newline{}
\courseList{Foundations of Computer Systems, Embedded System Software Engineering, Machine Learning Signal Processing, Distributed Systems, Autonomous Control System, Storage Systems, Modern Computer Architecture and Design, Large-scale Distributed Machine Learning and Optimization, Introduction to Embedded Systems, Multi-Robot Planning and Coordination, How to Write Fast Code II, Software Development for Computational Science and Engineering}

\resumeEducationSubheading
{Tianjin University} {Tianjin, China}
{Bachelor of Science in Computer Science and Technology; \quad GPA: 3.84/4.0} {June 2023}
\newline{}
\courseList{Data Structure, Operating System, Computer Networks, Principles of Database, Parallel Computing, Natural Language Processing, Pattern Recognition and Deep Learning}

\resumeEducationSubHeadingListEnd

\vspace{-9pt}

\section*{Publications}
\resumePublicationsSubHeadingListStart
\resumePublicationSubItem{Zhang, Z., \textbf{Guo, S.}, et al.}{``Bio-inspired Distributed Neural Locomotion Controller (D-NLC) for Robust Locomotion and Emergent Behaviors'' (May '25). \textit{2025 IEEE International Conference on Robotics and Automation} (ICRA 2025).}
\vspace{-10pt}
\resumePublicationSubItem{Bingdao, F., Di, J., et al.}{``Backdoor attacks on unsupervised graph representation learning'' (August '24). \textit{Journal of Neural Networks} (Vol. 180, p.106668).}
\resumePublicationSubItem{Jin, D., Feng, B., \textbf{Guo, S.}, et al.}{``Local-Global Defense Against Unsupervised Adversarial Attacks on Graphs'' (June '23). In the \textit{Thirty-Seventh AAAI Conference on Artificial Intelligence} (Vol. 37, No. 7).}
\resumePublicationSubItem{Ding, Y., \& \textbf{Guo, S.}}{``Conditional generative adversarial networks: Introduction and application'' (November '22). In the 2nd \textit{International Conference on Artificial Intelligence, Automation, and High-Performance Computing} (AIAHPC 2022) (Vol. 12348, pp. 258-266). SPIE.}
\resumePublicationsSubHeadingListEnd

\vspace{4pt}

%---------- RESEARCH & WORK EXPERIENCE -----------
\section{Research \& Work Experience}
\resumeResearchSubHeadingListStart

% ---- Yondu AI (current) ----
\resumeResearchSubheadingShort
{Yondu AI \textnormal{\small{$|$ Y Combinator, Early-stage Humanoid Robotics Startup}}} {Los Angeles, CA}
{Robot Intern (Full-time)} {Feb 2026 - Present}
\resumeResearchMidItemListStart
\resumeResearchMidItem{}
{Parallelized Motion Planning for Wheeled Humanoid Warehouse Robots}
{
  \resumeResearchItemListStart
  \resumeResearchItem{}
  {Developing parallelized \textbf{C++ inverse kinematics (IK) and motion planning pipelines} for wheeled humanoid warehouse robots, targeting real-time performance for autonomous manipulation tasks.}
  % \resumeResearchItem{}
  % {Architecting \textbf{ROS-based} system infrastructure integrating multi-threaded IK solvers with task-and-motion planning modules for warehouse operation.}
  % \resumeResearchItem{}
  {Designed \textbf{experience-based roadmap cache} system with persistence to accelerate similar pose/trajectory queries, reducing IK solving time for repeated warehouse tasks.}
  \resumeResearchItem{}
  {Developed a warehouse cart scanning simulation pipeline with QR code         
  localization, IK dispatch for barcode poses, and live virtual camera feedback during scanning.}
  \resumeResearchItemListEnd
}
\resumeResearchMidItem{}
{Ouster LiDAR Integration \& Dynamic IMU Estimation and Calibration System}
{
  \resumeResearchItemListStart
  \resumeResearchItem{}
  {Resolved persistent \textbf{asymmetric yaw drift and odometry snaping} in LiDAR-inertial SLAM through: \textbf{extrinsic calibration} (KISS-ICP), \textbf{EKF sensor fusion audit}, \textbf{dynamic IMU bias modeling and correction}, and \textbf{LiDAR point cloud deskewing}.}
  \resumeResearchItemListEnd
}
\resumeResearchMidItem{}
{Data Collection and Post-processing Pipeline Optimization \& Acceleration}
{
  \resumeResearchItemListStart
  \resumeResearchItem{} % optimized NVENC encoding, and Numba parallel reduction,
  {Redesigned cloud-to-dataset post-processing pipeline achieving \textbf{5-7x speedup} through CPU/GPU parallelism, breaking through computational bottlenecks and improving throughput from an average of \textbf{28 to 145 FPS}.}
  \resumeResearchItem{}
  {Implemented \textbf{parallel segment download} saturating ISP downlink and minimum-coverage data selection for robot training
  datasets, reducing cloud storage transfer by 92\%.}
  \resumeResearchItemListEnd
}
\resumeResearchMidItemListEnd

% ---- Cepton ----
\resumeResearchSubheadingShort
{Cepton, Inc. \textnormal{\small{$|$ LiDAR Technology Company}}} {San Jose, CA}
{Embedded Software Engineer (promoted from Software Engineer Intern, Full-time)} {May 2025 - Jan 2026}
\resumeResearchMidItemListStart
\resumeResearchMidItem{}
{MST (Magnosteer Control) System Development}
{
  \resumeResearchItemListStart
  \resumeResearchItem{}
  {Developed embedded \textbf{PD control with feedforward} for precision mirror movement control in LiDAR systems, implementing encoder auto-calibration and trajectory interpolation in \textbf{C/C++}.}
  \resumeResearchItem{}
  {Optimized embedded system architecture by redesigning \textbf{SPI transfers and event queues}, achieving \textbf{73\% CPU utilization reduction} (90\% $\rightarrow$ 24\%) on ARM-based controllers.}
  \resumeResearchItem{}
  {Built comprehensive \textbf{control simulation framework} in C and Python for algorithm validation, enabling systematic benchmarking of PD vs. PD-with-feedforward control strategies.}
  \resumeResearchItem{}
  {Engineered real-time \textbf{MST GUI and debug logging system} in \textbf{Rust}, providing field engineers with live system monitoring and parameter tuning capabilities.}
  \resumeResearchItem{}
  {Delivered \textbf{72\% power consumption reduction} (2.9W $\rightarrow$ 0.8W) on Ultra B3.3 120$^\circ$ LiDAR units through embedded firmware optimization and Python-based power analysis tools.}
  \resumeResearchItemListEnd
}
\resumeResearchMidItemListEnd

% ---- EigenBot ----
\resumeResearchSubheading
{Eigenbot - Bio-inspired Distributed Control for Modular Robot} {Pittsburgh, PA}
{Biorobotics Lab Research Assistant (Part-time)} {Dec 2023 - Oct 2024}
{Advisor: Professor Howie Choset, Research Scientist Lu Li} {Carnegie Mellon University}
\resumeResearchMidItemListStart
\resumeResearchMidItem{}
{Modular Robot Bio-inspired Curve Walking Implementation}
{
  \resumeResearchItemListStart
  % \resumeResearchItem{}
  % {Implemented the curve walking algorithm on the Eigenbot (a hexapod robot) in CoppeliaSim, inspired by six-leg insects' behavior. }
  % \resumeResearchItem{}
  % {Fine-tuned the weights of my algorithm to achieve a smooth curve walking gait, and integrated curve walking / steering into Eigenbot's behavior tree.}
    \resumeResearchItem{}
  {  Designed and implemented a bio-inspired curve walking algorithm in
  CoppeliaSim for hexapod locomotion,
  optimizing controller weights for smooth gait transitions and
  adding steering control into Eigenbot's behavior tree.}

  \resumeResearchItemListEnd
}
\resumeResearchMidItem{}
{Eigenbot Firmware Development and Testing}
{
  \resumeResearchItemListStart
  \resumeResearchItem{}
  {Developed Eigenbot EEPROM parameters auto-setting tools, which accelerated Hexapod gaits tuning process.}
  \resumeResearchItem{}
  {Tested the firmware. Redesigned and debugged the module communication protocol to ensure the robot's stability and reliability. Conducted the experiments with Logan, analyzed the gaits and fine-tuned firmware parameters.}
  \resumeResearchItem{}
  {Integrated Foot Sensor Feedback into the robot's message protocol.}
  \resumeResearchItem{}
  {Collected EigenBot's data from the real-world using OptiTrack, developing a pipeline to assess metrics like Gait Phase Difference and Leg Stance Duration, benchmarking its performance with centralized predefined gait.}
  \resumeResearchItemListEnd
}
\resumeResearchMidItemListEnd

% % ---- MMPUG ----
% \resumeResearchSubheading
% {MMPUG - Multi-Model Perception Uber Good} {Pittsburgh, PA}
% {Biorobotics Lab Research Assistant (Part-time)} {May 2024 - Sep 2024}
% {Advisor: Dr. Matthew J. Travers} {Carnegie Mellon University}
% \resumeResearchMidItemListStart
% \vspace{-5pt}
% \resumeResearchMidItem{}
% {FAR-Planner Implementation in a Heterogeneous Agent System with RC cars and legged robots}
% {
%   \resumeResearchItemListStart
%   \resumeResearchItem{}
%   {Implemented A* \& Theta* global planner with 2.5D orientation optimization for MMPUG architecture.}
%   \resumeResearchItem{}
%   {Integrated features to make the planner adaptive to dynamic obstacles and ensure a safety distance for the robots.}
%   \resumeResearchItem{}
%   {Deployed this Global Planner on the RC cars, and tested the planner both in simulation and in the real world.}
%   \resumeResearchItemListEnd
% }
% \resumeResearchMidItemListEnd

% ---- GNN / TJU ----
\resumeResearchSubheading
{Theory and Applications of Graph Convolutional Network} {Tianjin, China}
{Research Assistant (Part-time)} {Nov 2021 - May 2023}
{Advisor: Associate Professor Jin, Di} {Tianjin University}
\resumeResearchItemListStart
\vspace{-5pt}
\resumeResearchItem{}
{Investigated the theoretical relationship between GAT and GCN within Deep Graph Infomax (DGI); integrated mutual information maximization and unsupervised learning to improve graph representation quality.}
\resumeResearchItem{}
{Benchmarked adversarial robustness of GNNs using Metattack and Nettack attack matrices; reproduced and compared defense methods including RoSA, Pro-GNN, PA-GNN, GraphCL, and GraphSS.}
\resumeResearchItem{}
{Designed a novel \textbf{unsupervised backdoor attack} on GNNs based on GIN and Momentum Contrast, formulated the threat model, derived the loss function, and implemented trigger injection; achieved higher Attack Success Rate and lower Clean Accuracy Drop than prior unsupervised graph attacks. % Published in \textit{Journal of Neural Networks} (2024) and \textit{AAAI 2023}.
}
\resumeResearchItemListEnd

\resumeResearchSubHeadingListEnd

\vspace{-6pt}

%---------------- PROJECTS -----------------
\section{Projects}

\resumeProjectsSubHeadingListStart
\resumeProjectsSubHeading{{CloudFS, a Deduplicated Cloud File System} {| \href{https://github.com/Guo-lab/CloudFS_Design}{CMU 18-746 Course Project \faIcon{github}}}}{Oct 2024 - Dec 2024}
\resumeProjectsItemListStart
\resumeProjectsItem{Implemented CloudFS, leveraging the properties of local SSD and cloud storage to make data-placement decisions.}
\resumeProjectsItem{Identified and eliminated duplicate data with Rabin fingerprinting. Minimized storage usage and optimized data transfer in the form of segments. Designed buffer file to manually handle the file data transfer.}
\resumeProjectsItem{Implemented the IOCTL functions to support snapshot operations, such as create, restore, install, and uninstall.}
\resumeProjectsItem{Added in-memory LRU caching to further improve performance and reduce cloud operation costs, reaching \textbf{NO.1} on the scoreboard among all students.}
\resumeProjectsItemListEnd

\resumeProjectsSubHeading{{BusTub, a Relational DBMS} {| \href{https://github.com/Guo-lab/15645-Database}{Personal Project based on CMU 15-645 \faIcon{github}}}}{May 2024 - July 2024}
\resumeProjectsItemListStart
% \resumeProjectsItem{Implemented LRU-K policy, a disk scheduler and a buffer pool in the storage manager.}
% \resumeProjectsItem{Implemented disk-backed hash index in the DB system, using a variant of extendible hashing as the hashing scheme.}
% \resumeProjectsItem{Created the operator executors that execute SQL queries and implemented optimizer rules to transform query plans.}
% \vspace{-10pt}
% \resumeProjectsItem{Added transaction support by implementing optimistic multi-version concurrency control.}
\resumeProjectsItem{Implemented BusTub DBMS storage layer (LRU-K, disk scheduler, buffer   
    pool, extendible hash index), query execution (operator executors,     
    optimizer rules), and transaction support with optimistic multi-version
     concurrency control. }
\resumeProjectsItemListEnd

\vspace{-3pt}

\resumeProjectsSubHeading{{Distributed File System Design} {| \href{https://github.com/Guo-lab/15640-Distributed-Systems}{CMU 15-640 Course Project \faIcon{github}}}}{Feb 2024 - Apr 2024}
\resumeProjectsItemListStart
\resumeProjectsItem{Built a File-Caching Proxy over RMI with open-close session semantics, LRU cache replacement, and version-based freshness control; implemented low-level RPC serialization for NFS operations.}
\resumeProjectsItem{Designed a Two-phase Commit Protocol with failure recovery and lost-message handling; tuned a 3-tier auto-scaling web service and benchmarked system bottlenecks under dynamic load.}
\resumeProjectsItemListEnd

\resumeProjectsSubHeading{Blind Source Signal Decomposition \& Analysis | \href{https://github.com/Guo-lab/MLSP-G2}{CMU 18797 Research Project \faIcon{github}}}{Oct 2023 - Dec 2023}
\resumeProjectsItemListStart
\resumeProjectsItem{Filtered the mixed HDEMG signals and detected peaks to slice the peak-wave windows. Used the PCA algorithm to reduce the dimensions of the data and K-Means to cluster the potentially same Motor Units' signals.}
\resumeProjectsItem{Based on the clustered results, applied DTW, Covariance, Cosine in a multi-threshold way to compare the wave similarity. Verified the firing instances for each Motor Unit based on the comparison results.}
\resumeProjectsItemListEnd

\resumeProjectsSubHeadingListEnd

\vspace{-9pt}

%---------------- SKILLS -----------------
\section{Skills}
\vspace{-16pt}
\begin{multicols}{2}
  \begin{description}[itemsep=0pt]
    \resumeSkillsItem{Programming}{C/C++, Python, Rust, LaTeX, Verilog}
    \resumeSkillsItem{Frameworks}{ROS, TensorFlow/PyTorch, Isaac Sim, Drake}
    % \columnbreak
    \resumeSkillsItem{Tools}{GIT, Docker, MySQL, VSCode, Gem5, PSoC, Vitis}
    \resumeSkillsItem{Platforms}{MacOS, Linux, Arduino, GCP}
  \end{description}
\end{multicols}

\end{document}
